\documentclass[12pt]{article}
\usepackage[utf8]{inputenc}
\usepackage[T1]{fontenc}
\usepackage{hyperref}
\title{Parauk Wa Language Resources for K-12 Curriculum Development}
\author{Compiled by Claude Sonnet 5.5 (claude-sonnet-5-5) for Andr\'as Kornai}
\date{}
\begin{document}
\maketitle
\section{Introduction}
Parauk Wa is an Austroasiatic language of Myanmar and China, written in a Latin-based orthography devised by missionaries. A dictionary, a Bible and orthography notes were found; no specific primer was found.
\section{Primary Educational Resources}
The orthography page states that the revised Bible orthography was developed into literacy primers used in Wa Special Region 2, but no primer title or URL was seen.
\section{Linguistic and Reference Materials}
\begin{itemize}
\item Watkins, J., \emph{Dictionary of Wa} (2 vols., Brill, 2013): \url{https://eprints.soas.ac.uk/17317}
\item Wa orthography: \url{https://humancomp.org/wadict/wa_orthography.html}
\item Wa resources: \url{https://humancomp.org/wadict/wa_resources.html}
\item Wa language (Wikipedia): \url{https://en.wikipedia.org/wiki/Wa_language}
\end{itemize}
\section{Religious and Community Resources}
Parauk 2012 Edition Bible (Bible Society of Myanmar): \url{https://dbs.org/bibles/PRKBSM}
\section{Digital and Audio Resources}
Writing overview: \url{https://omniglot.com/writing/wa.htm}
\section{Curriculum Development Context}
The language is widely spoken, with an established Latin orthography and a major dictionary. Locally produced primers likely exist but were not located in these searches.
\end{document}